% COLLECTION_METADATA: {"schema_version": 1, "kind": "reused_writeup", "independently_reviewed": false, "primary_source": "attacks/erdos/gpt_pro_5.4/786.tex", "source_paths": ["attacks/erdos/gpt_pro_5.2/786.tex", "attacks/erdos/gpt_pro_5.4/786.tex"], "source_model": "gpt_pro_5.4", "source_completion": 100.0, "source_claim": "solved"}
\section*{Erd\H{o}s Problem 786 --- collection record}

\subsection*{PROVENANCE AND STATUS}
Official problem: https://www.erdosproblems.com/786

Saved collaborative-database status: open.
Snapshot checked: 2026-09-23T17:22:24Z.
Latest status and formalization links: https://teorth.github.io/erdosproblems/

This is an attributed collection of existing writeups. The model-folder names below identify their original attribution; placement in GPT\_6\_Astra\_Ultra does not attribute their mathematical work to that model. All locally available source versions selected for this problem are reproduced below, without editing their text.

Proof claims, computations, literature statements, and completion estimates in the reused text have not been independently verified in this collection pass. A source claiming a solution does not change the saved upstream status.

Primary source (latest numbered version in the preferred model folder): \texttt{attacks/erdos/gpt\_pro\_5.4/786.tex}.

\subsection*{REUSED SOURCE 1}
Original attribution: \texttt{gpt\_pro\_5.2}.
Original file: \texttt{attacks/erdos/gpt\_pro\_5.2/786.tex}.

% BEGIN REUSED SOURCE: attacks/erdos/gpt_pro_5.2/786.tex
\section*{Problem 786}
\addcontentsline{toc}{section}{Problem 786}

\subsection*{1. Formal restatement}
Let $A\subset \mathbb{N}$. We say that $A$ has the \emph{length-rigidity} property if for all integers $r,s\ge 1$ and all (not necessarily distinct) choices $a_1,\dots,a_r\in A$ and $b_1,\dots,b_s\in A$, the equality
\[
 a_1\cdots a_r=b_1\cdots b_s
\]
forces $r=s$.

\medskip
\noindent\textbf{(a) Infinite version.} Does there exist, for every $\varepsilon>0$, a set $A\subset\mathbb{N}$ with this property and natural density $d(A)>1-\varepsilon$?

\medskip
\noindent\textbf{(b) Finite version.} Let $A\subset\{1,2,\dots,N\}$ have the length-rigidity property. Must there be an absolute $c>0$ such that $|A|\le (1-c)N$ for all $N$?

\subsection*{2. Quick literature/context check}
The problem statement records the following constructions/bounds.
\begin{itemize}[leftmargin=2em]
\item The residue class $A=\{n\in\mathbb{N}: n\equiv 2\pmod 4\}$ has density $1/4$ and is length-rigid.
\item (Selfridge) There exist length-rigid sets of density $1/e-\varepsilon$ for any $\varepsilon>0$.
\item In the finite setting, if $A$ is the set of $n\le N$ having a prime factor $>\sqrt{N}$, then $|A|=(\log 2+o(1))N$ and $A$ is length-rigid.
\item (Ruzsa, unpublished as quoted in the statement) In the finite setting one should have $|A|\le (1-c)N$ for some absolute $c>0$.
\end{itemize}

One can further note (from later discussion in the literature) that taking $A$ to be integers $n\le N$ with \emph{exactly one} prime factor exceeding a suitably chosen threshold gives a larger constant-density construction; the best constant of this form is related to the Hall--Montgomery constant. (No attempt is made here to rederive that constant from first principles.)

\subsection*{3. Attack plan}
\begin{enumerate}[leftmargin=2em]
\item \textbf{Structural step.} Reinterpret length-rigidity as the existence of a completely additive ``length'' functional on prime-exponent vectors. This converts the multiplicative condition into linear algebra over $\mathbb{Q}$.
\item \textbf{Finite optimization heuristic.} For $A\subset[1,N]$, natural candidates are level sets of the form
\[
F(n)=\sum_{p\le N} \alpha_p v_p(n)\equiv 1,
\]
where $v_p$ is the $p$-adic valuation and only primes in a certain range have nonzero weight. One then tries to maximize $|\{n\le N: F(n)=1\}|$ using sieve/asymptotic estimates.
\item \textbf{Upper bounds.} If one could show that any such level set has density bounded away from $1$ (uniformly in the weights), that would settle (a) negatively and imply the existence of $c>0$ in (b). A possible route is to combine the structural lemma with distribution results for additive functions (or with ``large prime factor'' decompositions).
\item \textbf{Computational experiments.} For small $N$, brute-force search can be done using the structural lemma (solve a rational linear system) instead of enumerating all multiplicative relations.
\end{enumerate}

\subsection*{4. Work}
\subsubsection*{4.1. A linear-algebra characterization (finite sets)}
For an integer $n$ and a prime $p$ write $v_p(n)$ for the exponent of $p$ in the prime factorization of $n$.

\begin{lemma}[Additive functional criterion]
\label{lem:additive-functional}
Let $A\subset\mathbb{N}$ be finite, and let $P$ be the finite set of primes dividing at least one element of $A$. For $a\in A$ write its exponent vector
\[
\mathbf v(a)\coloneqq (v_p(a))_{p\in P}\in \mathbb{Z}_{\ge 0}^{P}.
\]
Then $A$ is length-rigid if and only if there exist rational numbers $(\alpha_p)_{p\in P}\in\mathbb{Q}^{P}$ such that
\begin{equation}
\label{eq:alpha-eq}
\sum_{p\in P} \alpha_p\, v_p(a)=1\qquad\text{for every }a\in A.
\end{equation}
Equivalently, there exists a completely additive function $F:\mathbb{N}\to\mathbb{Q}$ of the form
\[
F(n)=\sum_{p\in P} \alpha_p v_p(n)
\]
for which $F(a)=1$ for all $a\in A$.
\end{lemma}

\begin{proof}
First assume \eqref{eq:alpha-eq} holds. If $a_1\cdots a_r=b_1\cdots b_s$ with all $a_i,b_j\in A$, then comparing prime exponents gives
\[
\sum_{i=1}^r v_p(a_i)=\sum_{j=1}^s v_p(b_j)\qquad\text{for all }p\in P.
\]
Multiplying by $\alpha_p$ and summing over $p\in P$ yields
\[
\sum_{i=1}^r \sum_{p\in P} \alpha_p v_p(a_i)
=
\sum_{j=1}^s \sum_{p\in P} \alpha_p v_p(b_j).
\]
By \eqref{eq:alpha-eq} each inner sum equals $1$, so the left side is $r$ and the right side is $s$, hence $r=s$.

\medskip
Conversely, assume $A$ is length-rigid. Consider any rational linear relation among exponent vectors
\begin{equation}
\label{eq:relation}
\sum_{a\in A} c_a\,\mathbf v(a)=\mathbf 0,
\end{equation}
where $c_a\in\mathbb{Q}$. Let $D$ be a common denominator of the $c_a$ and set $m_a\coloneqq Dc_a\in\mathbb{Z}$. Then
$\sum_a m_a\mathbf v(a)=\mathbf 0$. Write $m_a=m_a^+-m_a^-$ with $m_a^+,m_a^-\in\mathbb{Z}_{\ge0}$.
The vector identity $\sum_a m_a\mathbf v(a)=\mathbf 0$ is equivalent to
\[
\prod_{a\in A} a^{m_a^+}=\prod_{a\in A} a^{m_a^-}.
\]
Both sides are products of elements of $A$ (allowing repetition). By length-rigidity the total number of factors on each side is equal:
\[
\sum_{a\in A} m_a^+=\sum_{a\in A} m_a^-.
\]
Hence $\sum_{a\in A} m_a=0$, and dividing by $D$ gives $\sum_{a\in A} c_a=0$.

Thus, any rational relation \eqref{eq:relation} forces $\sum_a c_a=0$, meaning that the assignment $\mathbf v(a)\mapsto 1$ is well-defined as a linear functional on the $\mathbb{Q}$-span of the vectors $\{\mathbf v(a):a\in A\}$. Therefore there exists a linear functional $L$ on $\mathbb{Q}^P$ with $L(\mathbf v(a))=1$ for all $a\in A$. Writing $L$ in coordinates gives coefficients $(\alpha_p)_{p\in P}$ satisfying \eqref{eq:alpha-eq}.
\end{proof}

\subsubsection*{4.2. Two concrete constructions (finite $[1,N]$)}
The lemma explains why ``one distinguished prime factor'' constructions work: one chooses weights $\alpha_p\in\{0,1\}$ and forces each $n\in A$ to have total weighted valuation $1$.

\begin{proposition}[Classical $\sqrt N$ threshold]
Fix $N\ge 1$ and set
\[
A_N\coloneqq\{n\in\{1,\dots,N\}: \text{$n$ has a prime factor $p>\sqrt N$}\}.
\]
Then $A_N$ is length-rigid.
\end{proposition}

\begin{proof}
If $p>\sqrt N$ divides some $n\le N$, then $p$ divides $n$ to exponent exactly $1$ and no other prime $>\sqrt N$ can divide $n$ (since two such primes would multiply to $>N$). Thus every $n\in A_N$ has
\[
\sum_{p>\sqrt N} v_p(n)=1.
\]
Apply Lemma~\ref{lem:additive-functional} with weights $\alpha_p=1$ for $p>\sqrt N$ and $\alpha_p=0$ otherwise.
\end{proof}

The size asymptotic $|A_N|=(\log 2+o(1))N$ comes from the classical estimate that the proportion of $\sqrt N$-smooth numbers up to $N$ tends to $1-\log 2$ (Dickman--de Bruijn theory); we do not reprove this here.

\subsubsection*{4.3. Small computational experiment (finite $[1,N]$)}
Using Lemma~\ref{lem:additive-functional}, one can test length-rigidity of a candidate $A\subset\{1,\dots,N\}$ by solving a rational linear system $\sum_p \alpha_p v_p(a)=1$ for all $a\in A$.

A brute-force search over all subsets for $N\le 12$ (implemented via exact rational Gaussian elimination) finds the following maximum sizes:
\[
\begin{array}{c|ccccccccccc}
N&2&3&4&5&6&7&8&9&10&11&12\\\hline
\max|A|&1&2&2&3&3&4&4&5&5&6&7
\end{array}
\]
For $N=12$ one extremal example is $A=\{3,5,6,7,10,11,12\}$; it is certified by the weights
$\alpha_3=\alpha_5=\alpha_7=\alpha_{11}=1$ and all other $\alpha_p=0$, i.e.
$F(n)=v_3(n)+v_5(n)+v_7(n)+v_{11}(n)$.

\subsection*{5. Verification}
\begin{itemize}[leftmargin=2em]
\item Lemma~\ref{lem:additive-functional} is internally consistent: it correctly rules out sets containing $1$ or simultaneously containing $p$ and $p^2$.
\item The construction $A_N$ is verified by the fact that each $n\in A_N$ has a unique prime factor $>\sqrt N$.
\item The small-$N$ brute force uses the lemma as a certificate, so any reported $A$ indeed satisfies the defining implication.
\end{itemize}

\subsection*{6. FINAL}
\textbf{UNRESOLVED.}

\medskip
\noindent\textbf{Fail-safe (required):}
\begin{enumerate}[leftmargin=2em,label=(\roman*)]
\item \textbf{Strongest partial results proved here.}
Lemma~\ref{lem:additive-functional} gives a clean equivalence between length-rigidity and the existence of a rational completely additive ``length'' functional constant on $A$. This gives a simple proof of the $\sqrt N$-threshold construction in the finite problem and supports systematic search.
\item \textbf{First gap that blocks a full solution.}
Even with the additive-functional characterization, we do not have a general theorem bounding the density of a level set $\{n: F(n)=1\}$ of a (possibly infinitely-supported) additive function $F(n)=\sum_p \alpha_p v_p(n)$.
\item \textbf{What would likely close the problem.}
A uniform upper bound of the form $d(\{n: F(n)=1\})\le 1-c$ for all nontrivial additive $F$ with rational values would settle (a) negatively and imply (b). Alternatively, an explicit construction achieving densities $1-\varepsilon$ would settle (a) positively.
\item \textbf{Why computations didn't settle it.}
The brute force only explores very small $N$ and cannot distinguish between densities such as $0.83$ vs $1$; it is useful mainly as a sanity check and for pattern-spotting.
\end{enumerate}

\subsection*{7. COMPLETION ESTIMATE}
\textbf{35\%.} A structural equivalence and supporting computations are provided, but the main density questions remain open.


% END REUSED SOURCE: attacks/erdos/gpt_pro_5.2/786.tex

\subsection*{REUSED SOURCE 2}
Original attribution: \texttt{gpt\_pro\_5.4}.
Original file: \texttt{attacks/erdos/gpt\_pro\_5.4/786.tex}.

% BEGIN REUSED SOURCE: attacks/erdos/gpt_pro_5.4/786.tex
\section{Problem 786}

\subsection*{1) ROUND-2 OBJECTIVE}
Path \textbf{(C)}: identify the fatal obstruction in Round 1 and correct the statement. In Round 1 I treated the repetitions-allowed version as still open. That is incorrect. The currently curated literature record shows that, when repetitions are allowed, the problem is already settled negatively: every such set is contained in a level set of a real-valued additive function, and Erd\H{o}s--Ruzsa--S\'ark\"ozy imply an absolute upper-density barrier. The minimal corrected open problem is the \emph{distinct-elements} version.

\subsection*{2) Round-1 FOUNDATION USED}
I use the following Round-1 results without reproving them.
\begin{enumerate}
\item The finite additive-functional criterion: for finite $A$, the repetitions-allowed property is equivalent to the existence of a rational additive functional on prime-exponent vectors taking the value $1$ on every $a \in A$.
\item The Round-1 diagnosis that the real gap was to control level sets of additive functions.
\item The standard examples (the class $n \equiv 2 \pmod 4$, Selfridge's construction, and the one-large-prime finite construction).
\end{enumerate}

\subsection*{3) NEW INSIGHT / TOOL (ROUND-2)}
The new ingredient is an \emph{infinite-set extension} of the Round-1 structural lemma. Instead of working only with finitely many primes and rational coefficients, I pass to the rational vector space of finitely supported prime-exponent vectors and extend a linear functional from the span of $\{\mathbf v(a): a \in A\}$ to the whole space. This gives a real-valued additive function on all natural numbers.

The second new ingredient is the corrected literature input: the updated 2026 discussion of Problem 786 records that Erd\H{o}s--Ruzsa--S\'ark\"ozy already settle the repetitions version negatively, and that the remaining open variant is the distinct-elements interpretation.

\subsection*{4) ATTACK PLAN (ROUND-2)}
The exact gap left by Round 1 was this: the finite linear-algebra lemma had not yet been upgraded to arbitrary $A$, so it could not be combined with the known distribution theory of additive functions. I will:
\begin{enumerate}
\item extend the structural lemma from finite $A$ to arbitrary $A$;
\item reduce the repetitions-allowed problem to a level set $\{n : f(n)=1\}$ of a real additive function $f$;
\item import the Erd\H{o}s--Ruzsa--S\'ark\"ozy consequence recorded in the current literature summary, obtaining a negative answer to both the density and finite problems;
\item isolate precisely why this argument does \emph{not} resolve the distinct-elements variant.
\end{enumerate}

\subsection*{5) WORK (ROUND-2)}
Write the repetitions-allowed property as follows:
\[
 a_1 \cdots a_r = b_1 \cdots b_s \quad (a_i,b_j \in A)
 \implies r=s.
\]

\paragraph{Lemma 786.1 (global additive-functional criterion).}
Let $A$ be a subset of the natural numbers with the repetitions-allowed property. Then there exists a real-valued additive function $f$ on the natural numbers such that
\[
 f(a)=1 \qquad \text{for every } a \in A.
\]

\paragraph{Proof.}
Let $P$ be the set of primes dividing at least one element of $A$. Let $V$ be the vector space over the rationals consisting of all finitely supported tuples
\[
 x=(x_p)_{p \in P}, \qquad x_p \in \mathbf Q.
\]
For a natural number $n$, let
\[
 \mathbf v(n) := (v_p(n))_{p \in P} \in V,
\]
where $v_p(n)$ is the exponent of $p$ in the factorisation of $n$.

Let $W$ be the rational span of $\{\mathbf v(a): a \in A\}$ inside $V$. Define a map $L$ on $W$ by declaring
\[
 L\!\left(\sum_{i=1}^m q_i \mathbf v(a_i)\right) := \sum_{i=1}^m q_i
\]
for rational numbers $q_i$ and elements $a_i \in A$.

I claim this is well-defined. Suppose
\[
 \sum_{i=1}^m q_i \mathbf v(a_i)=0.
\]
Choose a common denominator $D$ and set $m_i = D q_i \in \mathbf Z$. Then
\[
 \sum_{i=1}^m m_i \mathbf v(a_i)=0.
\]
Write $m_i = m_i^+ - m_i^-$ with $m_i^+,m_i^- \ge 0$. The equality of exponent vectors is equivalent to
\[
 \prod_{i=1}^m a_i^{m_i^+} = \prod_{i=1}^m a_i^{m_i^-}.
\]
Each side is a product of elements of $A$, with repetitions allowed. By the defining property of $A$, the total number of factors on the two sides must be equal:
\[
 \sum_{i=1}^m m_i^+ = \sum_{i=1}^m m_i^-.
\]
Hence
\[
 \sum_{i=1}^m m_i = 0,
\]
and dividing by $D$ gives $\sum_i q_i = 0$. So $L$ is well-defined.

Now extend $L$ from the subspace $W$ to a linear functional on all of $V$ by choosing a basis of $W$ and extending it to a basis of $V$. Keep the same name $L$ for the extension. Define
\[
 f(n) := L(\mathbf v(n)).
\]
Because
\[
 \mathbf v(mn)=\mathbf v(m)+\mathbf v(n),
\]
we have
\[
 f(mn)=f(m)+f(n)
\]
for all natural numbers $m,n$; that is, $f$ is additive. And for every $a \in A$,
\[
 f(a)=L(\mathbf v(a))=1.
\]
This proves the lemma.
\hfill$\square$

\paragraph{Corollary 786.2.}
Every repetitions-allowed set $A$ is contained in a level set of a real-valued additive function:
\[
 A \subseteq \{n : f(n)=1\}
\]
for some additive $f$.

\paragraph{External input.}
The current 2026 update to the Erd\H{o}s Problems page for Problem 786 records the following consequence of Erd\H{o}s--Ruzsa--S\'ark\"ozy (1973): once a set is contained in a fixed level set of a real-valued additive function, one obtains
\[
 \overline d(A) \le \frac12,
\]
and, in the finite problem, there is an absolute constant $c>0$ such that for all sufficiently large $N$,
\[
 |A \cap [1,N]| \le (1-c)N.
\]
Applying this to Corollary 786.2 gives the promised negative answers.

\paragraph{Theorem 786.3 (repetitions version is already settled).}
For the repetitions-allowed interpretation of Problem 786:
\begin{enumerate}
\item there is no set $A$ of density $>1-\varepsilon$ for every $\varepsilon>0$; in fact every such $A$ has upper density at most $1/2$;
\item in the finite problem one cannot have $|A|=(1-o(1))N$; indeed $|A| \le (1-c)N$ for some absolute $c>0$ and all sufficiently large $N$.
\end{enumerate}

\paragraph{Why the distinct-elements version is genuinely different.}
The proof above uses arbitrary integer relations among exponent vectors. Such a relation becomes a product identity only because repetitions are allowed. If one insists that all $a_i$ and all $b_j$ be distinct, then only disjoint $0$--$1$ relations are relevant; integer coefficients larger than $1$ are no longer legal. So the additive-function reduction collapses exactly at this point. This isolates the minimal corrected open statement: require the products on both sides to use distinct elements.

\subsection*{6) ADVERSARIAL VERIFICATION}
\begin{enumerate}
\item The definition of $L$ uses only finitely many elements of $A$ at a time, so no infinite product or convergence issue appears.
\item The extension step is safe because $W$ is a rational vector subspace of $V$; I am extending a linear functional on vector spaces, not claiming that a homomorphism to the integers extends on a free abelian group. This avoids the standard obstruction exemplified by $2\mathbf Z \to \mathbf Z$.
\item The conclusion is only $A \subseteq \{f=1\}$, not equality; this is the correct statement needed for the upper-bound transfer.
\item The distinct-elements obstruction is real: the argument from an integer relation
\[
 \sum m_a \mathbf v(a)=0
\]
produces repeated factors $a^{m_a^+}$ and $a^{m_a^-}$, which are illegal in the distinct-elements variant whenever some coefficient has absolute value bigger than $1$.
\end{enumerate}

\subsection*{7) FINAL}
\textbf{FULL SOLUTION -- COUNTEREXAMPLE / DISPROOF.}

For the repetitions-allowed interpretation, both optimistic conclusions are false. Every such set is contained in a level set of a real-valued additive function, and the Erd\H{o}s--Ruzsa--S\'ark\"ozy consequences recorded in the current literature give upper density at most $1/2$ and a finite upper bound $|A| \le (1-c)N$. The minimal corrected open problem is the distinct-elements version.

\subsection*{8) COMPLETION ESTIMATE}
\textbf{COMPLETION: 100\%}

\subsection*{9) REFERENCES}
\begin{thebibliography}{99}
\bibitem{ERS73-786}
P. Erd\H{o}s, I. Ruzsa, and A. S\'ark\"ozy,
\emph{On the number of solutions of $f(n)=a$ for additive functions},
Acta Arith. 24 (1973), 1--9.

\bibitem{Bloom786}
T. F. Bloom,
\emph{Erd\H{o}s Problem \#786}, updated 2 February 2026,
\texttt{https://www.erdosproblems.com/786}.

\bibitem{Bloom786Thread}
T. F. Bloom, T. Tao, and others,
\emph{Discussion thread for Erd\H{o}s Problem \#786}, especially comments of 18 October 2025 and 2 February 2026,
\texttt{https://www.erdosproblems.com/forum/thread/786}.
\end{thebibliography}


% END REUSED SOURCE: attacks/erdos/gpt_pro_5.4/786.tex

\subsection*{COMPLETION ESTIMATE}
Inherited primary-source estimate: $100\%$. This is the original source's unverified estimate, not a new assessment or confirmation that the problem is solved.
